\documentclass[../quyen]{subfiles}

\begin{document}

% ============================================================
% PHỤ LỤC B — ĐẶC TẢ USE CASE BỔ SUNG
% Tier 2 (9 UC compact spec) + Tier 3 (4 UC mention only)
% Tier 1 (11 UC full spec) đã ở Chương 2.3
% ============================================================

\section*{Đặc tả các Use case bổ sung}

Phụ lục này trình bày đặc tả của các use case Tier~2 (9 use case không quá phức tạp về mặt nghiệp vụ nhưng vẫn cần đặc tả chi tiết để phục vụ phát triển), và liệt kê ngắn gọn các use case Tier~3 (4 use case đơn giản --- mô tả 1 dòng cho đầy đủ).

\subsection*{B.1 Đặc tả UC Phê duyệt yêu cầu truy cập (UC012)}
\textit{\textbf{<TODO Stage H — bảng đặc tả UC012. Cite mobile/src/screens-v2/RequestsScreen.tsx handleApprove + EHRSystemSecure.sol approveRequestBySig>}}

\subsection*{B.2 Đặc tả UC Từ chối yêu cầu truy cập (UC013)}
\textit{\textbf{<TODO — bảng UC013 (Wave K sponsored). Cite RequestsScreen handleReject + rejectRequestBySig>}}

\subsection*{B.3 Đặc tả UC Đăng ký CCCD khẩn cấp (UC014)}
\textit{\textbf{<TODO — bảng UC014 + giải thích national ID hash. Cite backend/src/routes/profile.routes.js nationalIdSchema>}}

\subsection*{B.4 Đặc tả UC Thu hồi Người thân tin cậy (UC015)}
\textit{\textbf{<TODO — bảng UC015 + cascade revoke KeyShare source='trusted-contact'. Cite trustedContact.service.js removeContact>}}

\subsection*{B.5 Đặc tả UC Thu hồi uỷ quyền bác sĩ (UC016)}
\textit{\textbf{<TODO — bảng UC016 + epoch bump cascade. Cite DelegationScreen handleRevoke + ConsentLedger.sol revokeDelegation>}}

\subsection*{B.6 Đặc tả UC Tra cứu bệnh nhân khẩn cấp qua CCCD (UC017)}
\textit{\textbf{<TODO — bảng UC017. Cite mobile/src/screens-v2/doctor/EmergencyLookupScreen.tsx + backend/src/routes/emergency.routes.js>}}

\subsection*{B.7 Đặc tả UC Nộp chứng chỉ xác minh chuyên môn (UC018)}
\textit{\textbf{<TODO — bảng UC018. Cite mobile/src/screens-v2/doctor/CredentialSubmitScreen.tsx>}}

\subsection*{B.8 Đặc tả UC Xác minh bác sĩ độc lập (UC019)}
\textit{\textbf{<TODO — bảng UC019 (Wave E Ministry). Cite mobile/src/screens-v2/ministry/MinistryVerifyDoctorScreen.tsx + AccessControl.sol verifyDoctorByMinistry>}}

\subsection*{B.9 Đặc tả UC Tạm dừng / kích hoạt tổ chức (UC020)}
\textit{\textbf{<TODO — bảng UC020 (Wave F Ministry, typeword "THU HOI"). Cite MinistryDashboardScreen + setOrgActive>}}

\section*{Các Use case còn lại (Tier 3)}

\begin{itemize}
    \item \textbf{UC021 --- Loại bỏ thành viên bác sĩ khỏi tổ chức}: Quản trị viên Tổ chức thực hiện qua hàm \texttt{AccessControl.removeMember(orgId, doctor)}. Cite [mobile/src/screens-v2/org/OrgMembersScreen.tsx].
    \item \textbf{UC022 --- Thu hồi xác minh bác sĩ (Wave C)}: Quản trị viên Tổ chức gọi \texttt{revokeDoctorVerification(doctor)} → bác sĩ mất \texttt{VERIFIED\_DOCTOR} flag → \texttt{canAccess} refuse.
    \item \textbf{UC023 --- Thu hồi xác minh tổ chức (Wave F)}: Bộ Y tế gọi \texttt{revokeOrgVerification(orgId)} → org \texttt{verified} flag = false → mọi bác sĩ thuộc org mất \texttt{canAccess}.
    \item \textbf{UC024 --- Tạo phiên bản cập nhật hồ sơ của bệnh nhân (self)}: Bệnh nhân tự cập nhật, gọi \texttt{RecordRegistry.addRecord} với \texttt{parentCidHash} trỏ đến phiên bản trước. Tương đương UC001 nhưng có liên kết parent.
\end{itemize}

\end{document}
